\section{Экспериментальная часть}
\label{sec:experiments}

В настоящем разделе описана методика количественной оценки
разработанного решения, состав метрик, состав серии экспериментов
и полученные результаты.

\subsection{Методика проведения серии запусков}

Все запуски проводятся в одинаковой сцене и с одинаковой
конфигурацией Nav2/MPPI, различие между запусками определяется
случайными векторами скоростей подвижных препятствий и стохастикой
сэмплирования MPPI.

\subsection{Метрики}
\subsubsection{Распределение режимов сопровождения}

По данным топика \texttt{/target\_nav\_mode} строится временная диаграмма
режимов и интегральная доля времени каждого режима за запуск:
\begin{equation}
\bar{\tau}_m = \frac{1}{T} \int_{0}^{T} \mathbf{1}[m(t) = m]\,dt,
\quad m \in \{\text{Сопровождение}, \text{К~последней позиции},
\text{Поиск}\},
\label{eq:metric_modes}
\end{equation}
где $T$~--- длительность запуска. Высокая доля режима
\textit{Сопровождение цели} означает, что робот в большинстве
времени двигался непосредственно за пользователем, а не выполнял
поиск или перемещение к последней известной позиции.

\subsubsection{Видимость цели}

По бинарному сигналу $\nu(t) \in \{0,1\}$ из темы
\texttt{/target\_visible} вычисляется относительное время
наблюдаемости цели:
\begin{equation}
\bar{\nu} = \frac{1}{T} \int_{0}^{T} \nu(t)\,dt.
\label{eq:metric_visibility}
\end{equation}

\subsubsection{Метрики безопасности}

По данным топика \texttt{/scan} в каждый момент времени вычисляется
минимальное расстояние до препятствия
$d^{\text{lid}}_{\min}(t) = \min_i r_i(t)$. Введём два пороговых
расстояния: $d_{\text{col}}$~--- величина, ниже которой считается
произошедшим контактом с препятствием, и
$d_{\text{pot}} > d_{\text{col}}$~--- ширина зоны риска. По
последовательности значений $d^{\text{lid}}_{\min}(t)$ методом
подсчёта переходов с гистерезисом
\texttt{event\_release\_seconds} (исключение повторного срабатывания
на одном и том же приближении) определяются:
\begin{itemize}
    \item $N_{\text{c}}$~--- число эпизодов с $d^{\text{lid}}_{\min} \le d_{\text{col}}$
          (фактические контакты);
    \item $N_{\text{r}}$~--- число эпизодов входа в зону риска
          ($d^{\text{lid}}_{\min} \le d_{\text{pot}}$);
    \item $N_{\text{p}} = \max(0, N_{\text{r}} - N_{\text{c}})$~---
          число случаев, когда система зашла в зону риска, но
          избежала контакта;
    \item $\rho = N_{\text{p}}/N_{\text{r}}$ ($\rho := 1$ при
          $N_{\text{r}}=0$)~--- доля предотвращённых контактов.
\end{itemize}
Подсчёт ведётся начиная с первого момента, когда лидар фиксирует
$d^{\text{lid}}_{\min} > d_{\text{pot}}$, чтобы исключить артефакты
этапа спавна робота. В работе принято
$d_{\text{col}} = 0{,}30$~м,
$d_{\text{pot}} = 0{,}50$~м,
\texttt{event\_release\_seconds} $= 0{,}5$~с.

\subsubsection{Плавность управления}

По данным топика \texttt{/cmd\_vel} ($\mathbf{u}_k = [v_{x,k}, v_{y,k}, \omega_k]^{\top}$)
вычисляется норма приращения управления и её среднеквадратичное
значение по времени запуска:
\begin{equation}
\mathrm{RMS}(\|\Delta \mathbf{u}\|_2) =
\sqrt{\frac{1}{K-1}\sum_{k=1}^{K-1} \|\mathbf{u}_k - \mathbf{u}_{k-1}\|_2^{2}},
\label{eq:metric_smoothness}
\end{equation}
где $K$~--- число опубликованных управлений. Меньшие значения
соответствуют более плавному поведению контура управления и
непосредственно интерпретируются как минимизация слагаемого
$w_{\Delta u}\,\|\mathbf{u}_k - \mathbf{u}_{k-1}\|_2^{2}$ в~\eqref{eq:stage_cost}.

\subsection{Результаты одного запуска}

В качестве иллюстрации поведения системы на
рисунках~\ref{fig:exp_mode_run}--\ref{fig:exp_smooth_run} приведены
графики одного из запусков серии экспериментов.

\reportfig[1.00]{exp_mode_timeline_run.png}
{Временная диаграмма режимов сопровождения для отдельного запуска}
{fig:exp_mode_run}

На временной диаграмме режимов
(рисунок~\ref{fig:exp_mode_run}) видна типичная картина: робот в
основном находится в режиме \textit{Сопровождение цели} (зелёные
интервалы), кратковременно переходя в режим \textit{Поиск}, когда
цель временно исчезает из кадра, после чего возобновляет
сопровождение. Режим \textit{К последней позиции цели} активируется
эпизодически и кратко, что соответствует логике узла
\texttt{target\_nav\_bridge}.

\reportfig[1.00]{exp_visibility_run.png}
{Видимость цели $\nu(t)$ для отдельного запуска}
{fig:exp_vis_run}

График видимости (рисунок~\ref{fig:exp_vis_run}) показывает суммарную
долю времени с целью в кадре $\bar{\nu} \approx 69\%$, падения до
нуля соответствуют эпизодам, когда подвижные препятствия закрывали
цель, и совпадают с переключениями в режимы поиска на
рисунке~\ref{fig:exp_mode_run}.

\reportfig[1.00]{exp_safety_distance_run.png}
{Минимальное расстояние до препятствия по данным лидара
с горизонтальными порогами $d_{\text{pot}}$ и $d_{\text{col}}$}
{fig:exp_safety_run}

На графике безопасности (рисунок~\ref{fig:exp_safety_run}) видны
зоны кратковременного захода ниже $d_{\text{pot}}$, при этом для
данного запуска фактических контактов не зафиксировано
($N_{\text{c}}=0$, $N_{\text{r}}=3$, $N_{\text{p}}=3$, $\rho=1{,}00$).
Это иллюстрирует совместный эффект штрафов
$\Psi_{\text{obs},k}$~\eqref{eq:psi_obs} и
$\Psi_{\text{pf},k}$~\eqref{eq:psi_pf}: робот допускает заход в зону
риска, но не доводит ситуацию до контакта.

\reportfig[1.00]{exp_control_smoothness_run.png}
{Норма приращения управления $\|\Delta\mathbf{u}_k\|_2$ для отдельного запуска}
{fig:exp_smooth_run}

На графике плавности управления (рисунок~\ref{fig:exp_smooth_run})
повышенные значения в отдельных интервалах соответствуют моментам
маневрирования при обходе препятствий и переключений режимов,
длительные участки сопровождения дают близкие к нулю приращения
управления.

\subsection{Агрегированные результаты по серии экспериментов}

На рисунках~\ref{fig:exp_all_modes}--\ref{fig:exp_all_smooth}
приведены агрегированные характеристики по всем 40~запускам.

\reportfig[1.00]{exp_all_mode_shares.png}
{Доли времени по режимам сопровождения по запускам}
{fig:exp_all_modes}

Распределение режимов (рисунок~\ref{fig:exp_all_modes}) показывает,
что доминирующим режимом во всех запусках является
\textit{Сопровождение цели}, режимы \textit{Поиск} и
\textit{К~последней позиции цели} активируются только при
кратковременных потерях наблюдаемости и в большинстве запусков
занимают совокупно менее 25\% времени, что подтверждает корректность
логики переключения и устойчивость трекера.

\reportfig[1.00]{exp_all_visibility.png}
{Доля времени видимости цели $\bar{\nu}$ по запускам}
{fig:exp_all_vis}

Видимость цели (рисунок~\ref{fig:exp_all_vis}) в большинстве запусков
превышает 80\%, характерные просадки до 60--70\% соответствуют
сценариям, в которых подвижные препятствия систематически
закрывали цель, а единичная просадка ниже 20\% (\texttt{run32})~---
случаю, когда цель оказалась за углом помещения и удерживалась там
подвижным препятствием в течение значительной части запуска.

\reportfig[1.00]{exp_all_safety_events.png}
{Безопасность: число столкновений $N_{\text{c}}$,
заходов в зону риска $N_{\text{r}}$ и предотвращённых
контактов $N_{\text{p}}$ по запускам}
{fig:exp_all_safety}

Агрегированная характеристика безопасности
(рисунок~\ref{fig:exp_all_safety}) суммирует поведение системы по
всем 40~запускам. Видно, что:
\begin{itemize}
    \item число заходов в зону риска $N_{\text{r}}$ значимо превышает
          число фактических контактов $N_{\text{c}}$ практически в
          каждом запуске; это и есть прямой количественный эффект
          добавления направленного штрафа
          $\Psi_{\text{pf},k}$~\eqref{eq:psi_pf};
    \item суммарная по серии доля предотвращённых контактов
          составляет $\bar{\rho} \approx 0{,}70$.
\end{itemize}

\reportfig[1.00]{exp_all_control_smoothness.png}
{Среднеквадратичная норма приращения управления
$\mathrm{RMS}(\|\Delta\mathbf{u}_k\|_2)$ по запускам}
{fig:exp_all_smooth}

Среднеквадратичная норма приращения управления
(рисунок~\ref{fig:exp_all_smooth}) распределена в диапазоне
0{,}1--1{,}1, запуски с повышенными значениями соответствуют
сценариям с большим числом маневрирований.

\subsection{Обсуждение результатов}

Результаты экспериментальной серии подтверждают:
\begin{enumerate}
    \item функциональную корректность контура восприятия и
          переключения режимов: средняя доля видимости цели
          превышает 80\%, а доминирующий режим~---
          \textit{Сопровождение цели};
    \item полезность направленного штрафа $\Psi_{\text{pf},k}$:
          число фактических контактов $N_{\text{c}}$ систематически
          меньше числа заходов в зону риска $N_{\text{r}}$, что не
          обеспечивается одной лишь изотропной составляющей
          $\Psi_{\text{obs},k}$ при использовании голономной
          омни-платформы;
    \item приемлемую плавность управления: $\mathrm{RMS}$ приращений
          управления остаётся в пределах, обусловленных параметрами
          сэмплирования MPPI и ограничениями на ускорения.
\end{enumerate}